% from ega4-errata-addenda-list3-fr.tex
\oldpage[IV]{355}
\textup{(ii)} Compte tenu de
({\hyperref[canon-index.ega4.c04.thm.8-10-5]{8.10.5}}, (i)) appliqué suivant la
méthode de ({\hyperref[canon-index.ega4.c04.statement.8-1-2]{8.1.2}}, a)), on peut
se borner au cas où $S$ est un schéma local. Les réductions préalables
faites dans la preuve de (i) sont donc ici inutiles, et avec les mêmes
notations on voit que l'on a un diagramme commutatif
{\hyperlink{ega4.c04.sec.10-4-12-1}{(10.4.12.1)}} dans lequel $X$ est
un $S$-préschéma de type fini, $p$ est fini, $f'$ est un isomorphisme et
l'homomorphisme canonique $\sh{O}_X\to p_*(\sh{O}_{X'})$ est injectif.
Mais on montrera au chap.~VI, dans la théorie des quotients des
préschémas par des relations d'équivalence, que sous ces conditions $f$
est lui-même un isomorphisme (le lecteur vérifiera qu'on n'utilisera pas
{\hyperlink{ega4.c04.cor.10-4-12}{(10.4.12)}} dans la théorie des
relations d'équivalence du chap.~VI).
\end{proof}

\begin{corollary}[10.4.13]
\EGAReaderTarget{ega4.c04.cor.10-4-13}{IV.10.4.13}
Soient $k$ un corps, $X$ un préschéma algébrique sur $k$. Alors tout
$k$-endomorphisme radiciel $f$ de $X$ est fini et bijectif, donc un
homéomorphisme universel. Si de plus $X$ est réduit et $f$ birationnel
(par exemple si $k$ est de caractéristique 0), $f$ est un
$k$-automorphisme.
\end{corollary}

\begin{remark}[10.4.14]
\label{IV.10.4.14}
\begin{enumerate}[label=\textup{(\roman*)}]
\item Dans {\hyperlink{ega4.c04.cor.10-4-12}{(10.4.12)}} et
{\hyperlink{ega4.c04.cor.10-4-13}{(10.4.13)}} on ne peut supprimer
l'hypothèse que $X$ est réduit pour la seconde conclusion; c'est ce que
montre l'exemple du $B$-endomorphisme correspondant à l'augmentation
$D_B(L)\to B$ de l'anneau $D_B(L)=A$ défini dans ($0$, 18.2.3), où
$B$ est identifié à un sous-anneau de $A$.

\item Soit $X$ un préschéma noethérien. Nous ignorons si tout
endomorphisme radiciel $f$ de $X$ est nécessairement surjectif et fini.
Si $f$ est supposé surjectif et si de plus $X$ est un préschéma
japonais, le raisonnement de
{\hyperlink{ega4.c04.cor.10-4-12}{(10.4.12)}} montre que $f$ est fini;
si de plus $X$ est réduit et si $f$ est birationnel, par exemple si,
pour tout point maximal $x\in X$, $k(x)$ est parfait, le même
raisonnement montre que $f$ est un automorphisme.

\item On peut donner des exemples d'endomorphismes radiciels non
surjectifs de schémas affines intègres non noethériens : si $k$ est un
anneau, $A=k[(T_\lambda)_{\lambda\in L}]$ un anneau de polynômes où $L$
est infini dénombrable, et $\mathfrak J$ l'idéal engendré dans $A$ par
une sous-famille $(T_\lambda)_{\lambda\in H}$, où $L-H$ est infini et
$H\neq\varnothing$, on remarque que $A/\mathfrak J$ est isomorphe à
$k[(T_\lambda)_{\lambda\in L-H}]$, donc à $A$; l'endomorphisme de
$\operatorname{Spec}(A)$ correspondant à l'endomorphisme d'anneaux
composé
\[
 A\xrightarrow{p}A/\mathfrak J\xrightarrow{u}A
\]
où $p$ est la projection canonique et $u$ un isomorphisme, répond à la
question, son image étant le sous-préschéma fermé de
$\operatorname{Spec}(A)$ défini par $\mathfrak J$.
\end{enumerate}
\end{remark}

\paragraph{(Err\textsubscript{IV}, 32).}
Dans (IV, 11.3.14), ligne 7 de la p.~141, avant « Soient », insérer :
(i). Après la ligne 9 de la p.~141, ajouter :

\smallskip
\noindent
\textup{(ii)} Inversement, soit $f:X\to Y$ un morphisme plat. Si $X$
est unibranche (resp. géométriquement unibranche) en un point $x$, $Y$
est unibranche (resp. géométriquement unibranche) au point $y=f(x)$.

Au début de la ligne 10 de la p.~141, insérer : (i). Après la ligne 20
de la p.~141, ajouter :

\smallskip
\noindent
\textup{(ii)} Notons que si $Y_1=Y_{\mathrm{red}}$,
$X_1=X\times_Y Y_1$ a pour espace sous-jacent celui de $X$ et le
morphisme $f_1:X_1\to Y_1$ déduit de $f$ par changement de base est
plat. On peut donc se borner au cas où $Y=\operatorname{Spec}(A)$,
$A$ étant un anneau local réduit, et $y$ le point fermé de $Y$. Posons
$B=\sh{O}_{X,x}$; par hypothèse $B_{\mathrm{red}}$ est unibranche, donc
$\operatorname{Spec}(B)$ n'a qu'un seul point maximal, et par platitude
{\hyperref[canon-index.ega4.c04.prop.2-3-4]{(2.3.4)}} on en déduit qu'il en est de
même

